\documentclass[11pt]{amsart}
\usepackage[T1]{fontenc}
\usepackage{lmodern,microtype,amsmath,amssymb,amsthm,mathtools,enumitem}
\usepackage[margin=1.1in]{geometry}
\usepackage[hidelinks]{hyperref}
\usepackage{xurl}
\usepackage{booktabs,longtable}

\newtheorem{theorem}{Theorem}[section]
\newtheorem{lemma}[theorem]{Lemma}
\newtheorem{proposition}[theorem]{Proposition}
\newtheorem{corollary}[theorem]{Corollary}
\theoremstyle{remark}
\newtheorem{remark}[theorem]{Remark}
\numberwithin{equation}{section}

\newcommand{\EE}{\mathbb{E}}
\newcommand{\PP}{\mathbb{P}}
\newcommand{\RR}{\mathbb{R}}
\newcommand{\TW}{\mathrm{TW}_1}
\DeclareMathOperator{\Tr}{Tr}
\DeclareMathOperator{\Var}{Var}

\hypersetup{pdftitle={Tracy-Widom fluctuations for the row-centered conjugate kernel},pdfauthor={Yitzchak Shmalo}}

\title[Tracy--Widom law for the row-centered conjugate kernel]{Tracy--Widom fluctuations for the row-centered conjugate kernel of a random network with polynomial activation}
\author{Yitzchak Shmalo}
\address{Einstein Institute of Mathematics, The Hebrew University of Jerusalem, Jerusalem, Israel}
\email{yitzchak.shmalo@gmail.com}
\date{November 2026}
\subjclass[2020]{60B20, 15B52, 68T07}
\keywords{conjugate kernel, random features, sample covariance matrix, Tracy--Widom law, Gaussian hypercontractivity}

\begin{document}
\begin{abstract}
Let $Y=\sigma(WX/\sqrt d)$, where $W$ is an $m\times d$ matrix with independent standard Gaussian entries, the $N$ columns of $X$ are independent and uniform on the sphere of radius $\sqrt d$, and $\sigma$ is a fixed polynomial with $\EE\sigma(Z)=\EE[Z\sigma(Z)]=0$ for a standard Gaussian $Z$. Let $m,d,N\to\infty$ proportionally with $m/N\to b\ne1$, and let $K_P=YPY^{T}/N$, where $P$ removes the mean of each row of $Y$ over the data points. We prove that the largest eigenvalue of $K_P$ has Tracy--Widom fluctuations of order $N^{-2/3}$ around $v(1+\sqrt{m/N})^2$, where $v=\EE\sigma(Z)^2$, and that the same holds for $YY^{T}/N$ when $N/d$ times the squared coefficient of $\sigma$ on the normalized second Hermite polynomial $(z^2-1)/\sqrt2$ stays below $v\sqrt{N/m}$ in the limit. The proof applies, given the data, the edge universality theorem of Chen and Lyu with the quadratic form bound of Fan, Ma, Paquette and Wang, and shows that the covariance left after the row means are removed moves the right edge by $O(N^{-1})$.
\end{abstract}
\maketitle

\section{Introduction}\label{sec:intro}
The conjugate kernel of a one-layer network with random weights is the matrix of inner products of its hidden features. With weights $W\in\RR^{m\times d}$, data $X\in\RR^{d\times N}$ and activation $\sigma$ applied entrywise, the features form the $m\times N$ matrix $Y=\sigma(WX/\sqrt d)$, and the kernel is $K=YY^{T}/N$. When $m,d,N$ grow proportionally, its empirical spectral distribution converges. Pennington and Worah~\cite{PW} identified the limit for Gaussian weights and data; Louart, Liao and Couillet~\cite{LLC} gave deterministic equivalents for the resolvent, Benigni and P\'eché~\cite{BP21} treated general entry distributions, and Fan and Wang~\cite{FW} networks of any fixed depth. Benigni and P\'eché~\cite{BP22} proved, for $W$ and $X$ with independent identically distributed centered entries and an analytic activation, that the largest eigenvalue of the kernel has the same limit in probability as that of an information-plus-noise matrix; when the first Hermite coefficient of $\sigma$ vanishes the limit is explicit and has a phase transition that depends on $\sigma$ and on the kurtoses of the entries~\cite[Theorems~2.4 and~2.5]{BP22}. Wang, Wu and Fan~\cite{WWF} traced spiked input data to outlying eigenvalues of the kernel. Neither work gives the fluctuations of the largest eigenvalue. Here the data are uniform on the sphere. Theorem~\ref{thm:main} gives Tracy--Widom fluctuations for the largest eigenvalue of the row-centered kernel $K_P$ defined below, and Corollary~\ref{cor:K} gives them for $K$ when the largest eigenvalue of the covariance of a row stays below the threshold of Baik, Ben Arous and P\'ech\'e~\cite{BBP}.

Let $d,N,m$ be positive integers tending to infinity with
\begin{equation}\label{eq:ratios}
\frac Nd\to\alpha\in(0,\infty),\qquad b_N=\frac mN\to b\in(0,\infty)\setminus\{1\}.
\end{equation}
The columns $x_1,\dots,x_N$ of $X\in\RR^{d\times N}$ are independent and uniform on the sphere of radius $\sqrt d$, and $W\in\RR^{m\times d}$, independent of $X$, has independent standard Gaussian entries. Let $h_j=\mathrm{He}_j/\sqrt{j!}$ be the normalized Hermite polynomials, orthonormal in $L^2$ of the standard Gaussian measure. The activation is a fixed real polynomial
\begin{equation}\label{eq:sigma}
\sigma=\sum_{j=2}^{k}a_jh_j,\qquad v=\EE\sigma(Z)^2=\sum_{j=2}^{k}a_j^2>0,\qquad q=a_2^2,
\end{equation}
of degree $k\ge2$, with $Z$ standard Gaussian; equivalently $\EE\sigma(Z)=\EE[Z\sigma(Z)]=0$. The entries of $Y$ are $Y_{ij}=\sigma(w_i^{T}x_j/\sqrt d)$, where $w_i^{T}$ is the $i$-th row of $W$. Let $\mathbf 1$ be the vector of ones in $\RR^{N}$, $e=N^{-1/2}\mathbf 1$ and $P=I_N-ee^{T}$, so that $YP$ is $Y$ with the mean of each row subtracted, and let
\[
K_P=\frac1N\,YPY^{T}
\]
be the row-centered conjugate kernel.
We write $\TW$ for the Tracy--Widom distribution of the Gaussian orthogonal ensemble~\cite{TW}, $\lambda_{\max}$ and $\lambda_2$ for the largest and second largest eigenvalues of a symmetric matrix, $\|A\|_{\rm op}$ for the operator norm of a matrix $A$, and $\Longrightarrow$ for convergence in distribution.

Tracy--Widom fluctuations are known for sample covariance matrices whose sample vectors have dependent entries. Let $z_1,\dots,z_m\in\RR^{M}$ be independent and centered with covariance $T$, whose eigenvalues are $t_1\ge\dots\ge t_M>0$, and let $M/m$ tend to a limit other than one. The function $f(u)=-1/u+m^{-1}\sum_{j}t_j/(1+t_ju)$ on $(-1/t_1,0)$ of Silverstein and Choi~\cite{SC} has a derivative with one zero $u_+$; put $E_T=f(u_+)$, the right edge, and $\gamma_T=(2/f''(u_+))^{1/3}$ (Section~\ref{sec:edge}). Chen and Lyu~\cite[Corollary~1.4]{CL} prove that
\[
\gamma_Tm^{2/3}\Bigl[\lambda_{\max}\Bigl(\frac1m\sum_{i\le m}z_iz_i^{T}\Bigr)-E_T\Bigr]\Longrightarrow\TW
\]
when there is a fixed $\tau>0$ with $\tau\le t_M\le t_1\le\tau^{-1}$ and $1+t_1u_+\ge\tau$ (the right edge is then called regular), and when, for every $p\ge2$, there is $C_p$ with $\|z_i^{T}Az_i-\Tr(AT)\|_p\le C_p\|A\|_F$ for every deterministic complex symmetric $A$, where $\|V\|_p=(\EE|V|^{p})^{1/p}$ and $\|A\|_F$ is the Hilbert--Schmidt norm; we call this the quadratic form bound. Given the data, the rows of $Y$ are independent copies of $\sigma(U^{T}g)$, where $U=X/\sqrt d$ has the unit columns $u_j=x_j/\sqrt d$ and $g$ is a standard Gaussian vector in $\RR^{d}$; Fan, Ma, Paquette and Wang~\cite[Proposition~2.17]{FMPW} state the quadratic form bound for such vectors with high probability and a factor $N^{\epsilon}$, and their proof gives it with a constant independent of the dimension~\cite[proof of Proposition~2.17, by Lemma~5.7]{FMPW}; for Gaussian weights and polynomial $\sigma$, Proposition~\ref{prop:qf} gives a short proof with an explicit constant. Two things prevent a limit theorem for $K$. The covariance $\Sigma$ of a row given the data is $(v-q/d)I_N+(q/d)\mathbf 1\mathbf 1^{T}$ plus a matrix whose operator norm tends to zero in probability (Lemma~\ref{lem:cov}), so $\Sigma$ has an eigenvalue near $v+q\alpha$; when $q\alpha>v/\sqrt b$, the threshold of Baik, Ben Arous and P\'ech\'e~\cite{BBP}, the condition $1+t_1u_+\ge\tau$ fails. Below that threshold $E_T$ depends on the data, and a deterministic centering needs $E_T$ to within $o(N^{-2/3})$. Removing the row means removes the rank-one term: in an orthonormal basis of the vectors orthogonal to $\mathbf 1$, the rows of $YP$ have dimension $M=N-1$ and covariance $T=vI_M+A_0$ given the data, and $\|A_0\|_{\rm op}\to0$ in probability (Lemma~\ref{lem:T}). That convergence alone gives only $E_T-v(1+\sqrt{M/m})^2=o(1)$. The edge moves by an amount controlled by the trace and the Hilbert--Schmidt norm of $A_0$: for $M/m$ in a compact subset of $(0,\infty)$ there are $C,\epsilon_0>0$, depending only on $v$ and that set, such that $\|A_0\|_{\rm op}\le\epsilon_0$ implies (Lemma~\ref{lem:edge})
\[
\Bigl|E_T-v\bigl(1+\sqrt{M/m}\bigr)^2\Bigr|\le\frac CM\Bigl(|\Tr A_0|+\|A_0\|_F^2\Bigr).
\]
By Lemma~\ref{lem:T}, $\Tr A_0=-q(N-1)/d+o_P(1)$ and $\|A_0\|_F=O_P(1)$, where $o_P(1)$ tends to zero and $O_P(1)$ stays bounded in probability, so the right side is $O_P(N^{-1})$, below the scale $N^{-2/3}$ of the fluctuations.

\begin{theorem}\label{thm:main}
Let the columns of $X$ be independent and uniform on the sphere of radius $\sqrt d$, let $W$ be independent of $X$ with independent standard Gaussian entries, and assume~\eqref{eq:ratios} and~\eqref{eq:sigma}. Then
\begin{equation}\label{eq:main}
\frac{b_N^{1/6}}{v(1+\sqrt{b_N})^{4/3}}\,N^{2/3}\Bigl[\lambda_{\max}(K_P)-v\bigl(1+\sqrt{b_N}\bigr)^2\Bigr]\Longrightarrow\TW.
\end{equation}
In particular $\lambda_{\max}(K_P)\to v(1+\sqrt b)^2$ in probability.
\end{theorem}

The degree and the coefficients of $\sigma$ are fixed; there is no condition on the second Hermite coefficient $a_2$. Adding a constant to $\sigma$ adds a constant matrix to $Y$, which $P$ annihilates, so Theorem~\ref{thm:main} holds without the condition $\EE\sigma(Z)=0$, with $v=\Var\sigma(Z)$; only Corollary~\ref{cor:K} uses it. The centering uses $b_N$; replacing it by $b$ requires $N^{2/3}(b_N-b)\to0$.

\begin{corollary}\label{cor:K}
Under the hypotheses of Theorem~\ref{thm:main}, assume in addition $q\alpha<v/\sqrt b$. Then~\eqref{eq:main} holds with $K_P$ replaced by $K$.
\end{corollary}

\begin{remark}\label{rem:scope}
Since $K-K_P=\bar y\bar y^{T}$ with $\bar y=Y\mathbf 1/N$ is positive semidefinite of rank one, $\lambda_2(K)\le\lambda_{\max}(K_P)\le\lambda_{\max}(K)$, so Theorem~\ref{thm:main} also bounds $\lambda_2(K)$ from above at the scale $N^{-2/3}$ when $q\alpha\ge v/\sqrt b$.
\end{remark}

Given the data, $K$ and $K_P$ have, up to the factor $m/N$, the nonzero eigenvalues of sample covariance matrices $m^{-1}\sum_{i\le m}z_iz_i^{T}$ of $m$ independent vectors. For $z_i=T^{1/2}\xi_i$ with $T$ positive definite and independent entries in $\xi_i$, the largest eigenvalue of such a matrix has Tracy--Widom fluctuations: Johnstone~\cite{Johnstone} proved this for real Gaussian $\xi_i$ and $T=I$, with the law $\TW$ of Tracy and Widom~\cite{TW}. El Karoui~\cite{ElKaroui} proved it for complex Gaussian $\xi_i$ and a large class of $T$, with the unitary law of~\cite{TW94}, and Lee and Schnelli~\cite{LS} for real $\xi_i$ and general $T$. When $T$ is the identity except for a fixed number of eigenvalues below the threshold of Baik, Ben Arous and P\'ech\'e~\cite{BBP}, Bao, Pan and Zhou~\cite[Corollary~1.7]{BPZ} obtained $\TW$ for real entries, extending the diagonal case of F\'eral and P\'ech\'e~\cite{FP}; Corollary~\ref{cor:K} is a result of this kind for sample vectors with dependent entries.

Let the $M=N-1$ columns of a fixed matrix $Q$ be an orthonormal basis of the vectors orthogonal to $\mathbf 1$, so that $QQ^{T}=P$, and let $T=Q^{T}\Sigma Q$ be the covariance of $z=Q^{T}\phi$, $\phi=\sigma(U^{T}g)$. Since $z^{T}Az=\phi^{T}(QAQ^{T})\phi$, $\Tr(AT)=\Tr(QAQ^{T}\Sigma)$ and $\|QAQ^{T}\|_F=\|A\|_F$, a quadratic form bound for $\phi$ passes to $z$, and to complex symmetric $A$ by splitting $A$ into its real and imaginary parts. For Gaussian weights and polynomial $\sigma$, Proposition~\ref{prop:qf} gives one with an explicit constant: for $T=\EE zz^{T}$, every real symmetric $A$ and every $p\ge2$,
\[
\bigl\|z^{T}Az-\Tr(AT)\bigr\|_p\le2\,\|U\|_{\rm op}\,\bigl(\EE|\sigma'(Z)|^4\bigr)^{1/4}\,3^{k/2}\,\|T\|_{\rm op}^{1/2}\,(p-1)^{k}\,\|A\|_F ,
\]
and $\|U\|_{\rm op}$ and $\|T\|_{\rm op}$ are bounded with probability tending to one. Its proof bounds the variance of $z^{T}Az$ by the Gaussian Poincar\'e inequality, proved by Chernoff~\cite{Chernoff} in dimension one, and passes to the $p$-th moment by hypercontractivity of the Ornstein--Uhlenbeck semigroup, due to Nelson~\cite{Nelson} (see also Hariya~\cite{Hariya}). The hypothesis that $\sigma$ is a polynomial is used only in this proposition; Sections~\ref{sec:cov} and~\ref{sec:edge} use only the expansion of $\Sigma$ in powers of the inner products of the data. Ma and Misiakiewicz~\cite{MM} prove the anisotropic local law under the quadratic form bound alone.

Lemma~\ref{lem:edge} compares the function $f$ for $T$ and for $vI_M$, whose right edge $v(1+\sqrt{M/m})^2$ is that of the law of Marchenko and Pastur~\cite{MP}; the same lemma, with the largest eigenvalue of the covariance treated separately, covers $K$ below the threshold. If the Hermite coefficient $a_1=\EE[Z\sigma(Z)]$ of degree one is not zero, $\Sigma$ contains in addition the matrix $a_1^2U^{T}U$, whose spectrum is not concentrated near one point, and Lemma~\ref{lem:edge} does not apply.

\begin{center}
\small
\begin{tabular}{@{}lp{0.72\textwidth}@{}}\toprule
Section~\ref{sec:gauss} & hypercontractivity and the Gaussian Poincar\'e inequality for polynomials\\
Section~\ref{sec:qf} & the quadratic form bound for a deterministic $U$ (Proposition~\ref{prop:qf})\\
Section~\ref{sec:cov} & the covariance $\Sigma$ of a row given the data and the matrix $A_0$ (Lemmas~\ref{lem:cov} and~\ref{lem:T})\\
Section~\ref{sec:edge} & the displacement of the right edge (Lemma~\ref{lem:edge})\\
Section~\ref{sec:proof} & the proof of Theorem~\ref{thm:main} and Corollary~\ref{cor:K} on an event of probability tending to one, and the finite-$N$ edge (Remark~\ref{rem:finite})\\
Appendix~\ref{app:notation} & notation\\
\bottomrule
\end{tabular}
\end{center}

\section{Two Gaussian inequalities}\label{sec:gauss}
Let $g$ be a standard Gaussian vector in $\RR^{n}$ and $\|\cdot\|_p$ the norm of $L^p$.

\begin{lemma}\label{lem:hyper}
Let $H$ be a real polynomial of total degree at most $l$ in the coordinates of $g$. Then $\|H(g)\|_p\le(p-1)^{l/2}\|H(g)\|_2$ for every $p\ge2$.
\end{lemma}

\begin{proof}
Let $(S_t)$ be the Ornstein--Uhlenbeck semigroup. Nelson's hypercontractivity~\cite{Nelson}, in the form stated in the introduction of~\cite{Hariya}, gives $\|S_tG\|_p\le\|G\|_2$ for every square-integrable function $G$ of the Gaussian vector when $\exp(2t)=p-1$. Write $H=\sum_{j=0}^{l}H_j$, where $H_j$ is the component of $H$ in the $j$-th Wiener chaos, and put $G=\sum_{j=0}^{l}\exp(tj)H_j$. Since $S_tH_j=\exp(-tj)H_j$, we have $S_tG=H$, and by orthogonality of the chaoses $\|G\|_2^2=\sum_j\exp(2tj)\|H_j\|_2^2\le\exp(2tl)\|H\|_2^2=(p-1)^{l}\|H\|_2^2$.
\end{proof}

\begin{lemma}\label{lem:poincare}
For every real polynomial $H$ on $\RR^n$, $\Var H(g)\le\EE\|\nabla H(g)\|^2$.
\end{lemma}

\begin{proof}
This is the Gaussian Poincar\'e inequality~\cite{Chernoff}. For a polynomial on $\RR^n$ it can be read off the expansion $H=\sum_\nu c_\nu h_\nu$ in the orthonormal product Hermite basis, $h_\nu(g)=\prod_ih_{\nu_i}(g_i)$: the identity $\partial_ih_\nu=\sqrt{\nu_i}\,h_{\nu-\delta_i}$ gives $\EE\|\nabla H\|^2=\sum_\nu|\nu|\,c_\nu^2$, while $\Var H=\sum_{|\nu|\ge1}c_\nu^2$.
\end{proof}

\section{The quadratic form bound for Gaussian weights}\label{sec:qf}
In this section $\sigma$ is any real polynomial of degree $k$ with $\EE\sigma(Z)=0$. Let $U\in\RR^{d\times N}$ be deterministic with unit columns $u_1,\dots,u_N$ and $\|U\|_{\rm op}\le C_U$, and let $Q\in\RR^{N\times M}$ be deterministic with $Q^{T}Q=I_M$. For a standard Gaussian $g\in\RR^{d}$ let
\[
\phi(g)=\sigma(U^{T}g)\in\RR^N,\qquad z(g)=Q^{T}\phi(g)\in\RR^{M},\qquad T=\EE\,zz^{T},
\]
and assume $\|T\|_{\rm op}\le C_T$. Since each $u_j^{T}g$ is standard Gaussian, $\EE\phi=0$ and $\EE z=0$. Put $\kappa_4=(\EE|\sigma'(Z)|^4)^{1/4}$.

\begin{proposition}\label{prop:qf}
For every real symmetric $A\in\RR^{M\times M}$ and every $p\ge2$,
\begin{equation}\label{eq:qf}
\bigl\|z^{T}Az-\Tr(AT)\bigr\|_p\le2\,C_U\kappa_4\,3^{k/2}\sqrt{C_T}\,(p-1)^{k}\,\|A\|_F.
\end{equation}
For complex symmetric $A$ the same bound holds with an additional factor $\sqrt2$.
\end{proposition}

\begin{proof}
Let $F(g)=z(g)^{T}Az(g)$, so $\EE F=\Tr(AT)$, and let $L_j(g)=(QAz(g))_j$ for $j\le N$. By the chain rule
\[
\nabla F(g)=2\,U\operatorname{diag}\bigl(\sigma'(U^{T}g)\bigr)\,QAz(g).
\]
Lemma~\ref{lem:poincare}, the bound on $\|U\|_{\rm op}$ and the Cauchy--Schwarz inequality give
\begin{equation}\label{eq:var}
\Var F\le4C_U^2\sum_{j=1}^{N}\EE\bigl[\sigma'(u_j^{T}g)^2L_j(g)^2\bigr]\le4C_U^2\kappa_4^2\sum_{j=1}^{N}\|L_j\|_4^2,
\end{equation}
since each $\sigma'(u_j^{T}g)$ has the distribution of $\sigma'(Z)$. Each $L_j$ is a polynomial of degree at most $k$ in $g$, so Lemma~\ref{lem:hyper} with $p=4$ gives $\|L_j\|_4^2\le3^{k}\EE L_j^2$, and
\[
\sum_{j=1}^{N}\EE L_j^2=\EE\|QAz\|^2=\EE\,z^{T}A^2z=\Tr(A^2T)\le C_T\|A\|_F^2.
\]
Hence $\Var F\le4C_U^2\kappa_4^2\,3^{k}C_T\|A\|_F^2$. The polynomial $F-\EE F$ has degree at most $2k$, and Lemma~\ref{lem:hyper} gives $\|F-\EE F\|_p\le(p-1)^{k}(\Var F)^{1/2}$, which is~\eqref{eq:qf}. For complex symmetric $A=A_1+iA_2$ with $A_1,A_2$ real symmetric, apply~\eqref{eq:qf} to $A_1$ and $A_2$ and use $\|A_1\|_F+\|A_2\|_F\le\sqrt2\,\|A\|_F$.
\end{proof}


\section{The covariance of a row given the data}\label{sec:cov}
Write $\rho_{jl}=u_j^{T}u_l$ for the inner products of the columns of $U$. Given $U$, the rows of $Y$ are independent copies of $\phi(g)=\sigma(U^{T}g)$ with $g$ standard Gaussian in $\RR^{d}$. For standard Gaussian variables $Z_1,Z_2$ with correlation $\rho$, $\EE h_s(Z_1)h_t(Z_2)=\delta_{st}\rho^{s}$~\cite[Section~11.2]{ODonnell}, so the covariance $\Sigma=\EE[\phi\phi^{T}\mid U]$ of a row given the data is
\begin{equation}\label{eq:Sigma}
\Sigma_{jl}=\sum_{s=2}^{k}a_s^2\rho_{jl}^{s},\qquad\Sigma_{jj}=v.
\end{equation}
Let $B$ and $R$ be the symmetric matrices with zero diagonal and, for $j\ne l$,
\[
B_{jl}=\rho_{jl}^2-\frac1d,\qquad R_{jl}=\sum_{s=3}^{k}a_s^2\rho_{jl}^{s}.
\]
Then
\begin{equation}\label{eq:decomp}
\Sigma=\Bigl(v-\frac qd\Bigr)I_N+\frac qd\,\mathbf 1\mathbf 1^{T}+qB+R.
\end{equation}
The matrix $\Sigma$ is an inner-product kernel matrix of the data, and~\eqref{eq:decomp} is the splitting into an identity part, a rank-one part, a quadratic part and parts of higher degree that El Karoui~\cite{EK10} used to identify the spectrum of such matrices; Fan and Montanari~\cite{FM} bound the operator norm of the parts of higher degree. Lemma~\ref{lem:cov} bounds, with explicit constants, the Hilbert--Schmidt norm of $\Sigma-vI_N$, and Lemma~\ref{lem:T} turns this into bounds on $\Tr A_0$ and $\|A_0\|_F$, which govern the displacement in Lemma~\ref{lem:edge}; a bound on the operator norm does not control them.
For two independent uniform unit vectors in $\RR^{d}$, writing a standard Gaussian vector as its length times an independent uniform direction gives
\begin{equation}\label{eq:sphmom}
\EE\rho^{2s}=\frac{1\cdot3\cdots(2s-1)}{d(d+2)\cdots(d+2s-2)}.
\end{equation}

\begin{lemma}\label{lem:cov}
Let $D=\Sigma-vI_N$. Then $\Tr D=0$ and
\begin{align}
\EE\Tr D^2&\le\frac{3v^2N(N-1)}{d(d+2)},\label{eq:D2}\\
\EE\|R\|_F^2&\le\frac{15v^2N(N-1)}{d(d+2)(d+4)},\label{eq:R2}\\
\EE\Tr B^4&\le\frac{60N^2}{d^4}+\frac{8N^3}{d^4}+\frac{8N^4}{d^6}.\label{eq:B4}
\end{align}
In particular $\EE\Tr D^2=O(1)$, $\EE\|R\|_F^2=O(N^{-1})$ and $\|B\|_{\rm op}=O_P(N^{-1/4})$.
\end{lemma}

\begin{proof}
The diagonal of $D$ vanishes by~\eqref{eq:Sigma}. Off the diagonal $|D_{jl}|\le v\rho_{jl}^2$ and $|R_{jl}|\le v|\rho_{jl}|^3$, which with~\eqref{eq:sphmom} give~\eqref{eq:D2} and~\eqref{eq:R2}.

For~\eqref{eq:B4}, let $\mathcal S$ be the space of traceless symmetric $d\times d$ matrices with the inner product $\langle A,A'\rangle=\Tr(AA')$, and $\zeta(u)=uu^{T}-I_d/d\in\mathcal S$. Then $B_{jl}=\langle\zeta(u_j),\zeta(u_l)\rangle$ for $j\ne l$, and $\|\zeta(u)\|^2=c:=1-1/d$. By symmetry,
\[
\EE u_au_bu_cu_f=\frac{\delta_{ab}\delta_{cf}+\delta_{ac}\delta_{bf}+\delta_{af}\delta_{bc}}{d(d+2)},
\]
the constant being fixed by $\EE(\sum_au_a^2)^2=1$, so $\EE\langle\zeta(u),A\rangle^2=\EE(u^{T}Au)^2=2\Tr(A^2)/(d(d+2))$ for $A\in\mathcal S$: $\zeta(u)$ has mean zero and covariance $\theta I$ on $\mathcal S$, with $\theta=2/(d(d+2))$. In the expansion of $\Tr B^4$ over closed walks $j_1j_2j_3j_4$ with consecutive indices distinct, a walk on two vertices contributes $\EE B_{12}^4=\mu_4:=\EE(\rho^2-1/d)^4$, and there are $(N)_2$ of them. A walk on three vertices repeats one vertex; conditioning on that vertex gives $\EE B_{12}^2B_{13}^2=(c\theta)^2$, and there are $2(N)_3$ of them. A walk on four vertices contributes $\EE B_{12}B_{23}B_{34}B_{41}=\theta^2\EE B_{13}^2=c\theta^3$, integrating $u_2$ and $u_4$ first, and there are $(N)_4$ of them, where $(N)_i=N(N-1)\cdots(N-i+1)$. Thus
\[
\EE\Tr B^4=(N)_2\mu_4+2(N)_3(c\theta)^2+(N)_4c\theta^3.
\]
Expanding $\mu_4$ with~\eqref{eq:sphmom},
\[
\mu_4=\frac{12(d-1)(5d^2-13d+12)}{d^4(d+2)(d+4)(d+6)}\le\frac{60}{d^4}\qquad(d\ge2),
\]
because $(d-1)(5d^2-13d+12)=5d^3-18d^2+25d-12\le5d^3$ for $d\ge2$. With $c\le1$ and $\theta\le2/d^2$ this gives~\eqref{eq:B4}, and $\|B\|_{\rm op}^4\le\Tr B^4$.
\end{proof}

\begin{lemma}\label{lem:T}
Let $M=N-1$ and let $Q$ be the matrix of the introduction, so that $Q^{T}Q=I_M$, $QQ^{T}=P$ and $Q^{T}\mathbf 1=0$. The covariance of $z=Q^{T}\phi$ given $U$ is
\begin{equation}\label{eq:T}
T=Q^{T}\Sigma Q=\Bigl(v-\frac qd\Bigr)I_M+qQ^{T}BQ+Q^{T}RQ,
\end{equation}
and $A_0=T-vI_M=Q^{T}DQ$ satisfies
\begin{equation}\label{eq:A0}
\Tr A_0=-e^{T}De,\qquad\|A_0\|_F\le\|D\|_F,\qquad\|A_0\|_{\rm op}\le\frac qd+q\|B\|_{\rm op}+\|R\|_F,
\end{equation}
with $e^{T}De=q(N-1)/d+q\,e^{T}Be+e^{T}Re$, $|e^{T}Be|\le\|B\|_{\rm op}$ and $|e^{T}Re|\le\|R\|_F$. In particular $\|A_0\|_{\rm op}=o_P(1)$, $\Tr A_0=-q(N-1)/d+o_P(1)$ and $\|A_0\|_F=O_P(1)$.
\end{lemma}

\begin{proof}
Formula~\eqref{eq:T} is~\eqref{eq:decomp} with $Q^{T}\mathbf 1=0$. Since $\Tr D=0$, $\Tr A_0=\Tr(PD)=\Tr D-e^{T}De=-e^{T}De$, and $\|Q^{T}DQ\|_F\le\|D\|_F$ because $Q$ has orthonormal columns. The bound on $\|A_0\|_{\rm op}$ is read off~\eqref{eq:T}, using $\|R\|_{\rm op}\le\|R\|_F$. By~\eqref{eq:decomp}, $D=-(q/d)I_N+(q/d)\mathbf 1\mathbf 1^{T}+qB+R$, and $e^{T}e=1$, $(e^{T}\mathbf 1)^2=N$ give the formula for $e^{T}De$. The last statement follows from Lemma~\ref{lem:cov} and $N/d\to\alpha$.
\end{proof}

\section{Displacement of the right edge}\label{sec:edge}
Recall from Section~\ref{sec:intro} the edge function $f$ of a covariance $T$ of dimension $M$ with eigenvalues $t_1\ge\dots\ge t_M>0$, the zero $u_+$ of $f'$, the right edge $E_T=f(u_+)$, the scaling parameter $\gamma_T=(2/f''(u_+))^{1/3}$ and the constant $\tau$ of a regular right edge. Here
\[
f(u)=-\frac1u+\frac1m\sum_{j=1}^{M}\frac{t_j}{1+t_ju}
\]
on $(-1/t_1,0)$~\cite{SC,ElKaroui,LS}; this is the function of~\cite[(1.5)]{CL}, with $u$ in place of their $m$. There $f''(u)=-2/u^3+(2/m)\sum_jt_j^3/(1+t_ju)^3>0$, while $f'(u)\to-\infty$ as $u\downarrow-1/t_1$ and $f'(u)\to+\infty$ as $u\uparrow0$, so the zero $u_+$ of $f'$ is unique. Regularity is~\cite[Assumption~1.2]{CL}, with $\tau$ independent of $M$ and $m$. For $T=vI_M$ and $r=M/m$, $f$ becomes $f_0(u)=-1/u+rv/(1+vu)$, and
\[
u_0=-\frac1{v(1+\sqrt r)},\qquad E_0=v(1+\sqrt r)^2,\qquad\gamma_0=\frac{r^{1/6}}{v(1+\sqrt r)^{4/3}},\qquad f_0''(u_0)=\frac{2v^3(1+\sqrt r)^4}{\sqrt r},
\]
and for $v=1$ the inverse $\gamma_0^{-1}=(1+\sqrt r)(1+1/\sqrt r)^{1/3}$ is the scaling of Johnstone~\cite{Johnstone}; for $v=1$ these are the values in~\cite[Remark~1.5]{CL}.

\begin{lemma}\label{lem:edge}
Let $r=M/m$ lie in a fixed compact subset of $(0,\infty)$.
\begin{enumerate}
\item[(i)] Put $\epsilon=\max_j|t_j-v|$, $\mu=M^{-1}\sum_j(t_j-v)$ and $s_2=M^{-1}\sum_j(t_j-v)^2$. There are $\epsilon_0,C,c_0>0$, depending only on $v$ and the compact set, such that $\epsilon\le\epsilon_0$ implies
\begin{equation}\label{eq:edge}
|u_+-u_0|+|E_T-E_0|+|\gamma_T-\gamma_0|\le C(|\mu|+s_2),\qquad1+t_1u_+\ge c_0.
\end{equation}
\item[(ii)] Let $\eta\in(0,1)$, suppose $t_1\le(1-\eta)v(1+\sqrt r)$, and define $\epsilon$, $\mu$ and $s_2$ as in (i) with the sums and the maximum over $j\ge2$ only. There are $\epsilon_0,C,c_0>0$ and $m_0$, depending only on $v$, $\eta$ and the compact set, such that $\epsilon\le\epsilon_0$ and $m\ge m_0$ imply~\eqref{eq:edge} with $C(|\mu|+s_2)$ replaced by $C(|\mu|+s_2+m^{-1})$.
\end{enumerate}
\end{lemma}

\begin{proof}
(i) Choose a closed interval $J$ around $u_0$, uniformly over the compact range of $r$, such that $u<0$ on $J$, $1+vu\ge2c_0$ on $J$, $f_0''$ is bounded below on $J$ by a positive constant, and $f_0'$ takes opposite signs at the endpoints of $J$. For $\epsilon$ small, $1+tu\ge c_0$ on $J$ whenever $|t-v|\le\epsilon$, so $J\subset(-1/t_1,0)$. With $\delta=t-v$,
\[
\frac{t}{1+tu}-\frac{v}{1+vu}=\frac{\delta}{(1+vu)^2}-\frac{u\,\delta^2}{(1+vu)^2(1+tu)},
\]
and the last term and its first two derivatives in $u$ are bounded on $J$ by $C\delta^2$. Averaging over $j$ gives $\sup_J|f^{(l)}-f_0^{(l)}|\le C(|\mu|+s_2)$ for $l=0,1,2$. Since $|\mu|\le\epsilon$ and $s_2\le\epsilon^2$, after decreasing $\epsilon_0$ this is small enough that $f'$ takes opposite signs at the endpoints of $J$ and $f''$ is bounded below on $J$. Then $f'$ has a zero in $J$, which is $u_+$ by uniqueness, and the mean value theorem gives $|u_+-u_0|\le C(|\mu|+s_2)$. Writing $E_T-E_0=(f-f_0)(u_+)+f_0(u_+)-f_0(u_0)$ and $f''(u_+)-f_0''(u_0)=(f''-f_0'')(u_+)+f_0''(u_+)-f_0''(u_0)$, and using that $f_0'$ and $f_0'''$ are bounded on $J$, gives the bounds on $E_T-E_0$ and on $f''(u_+)-f_0''(u_0)$, and with them the bound on $\gamma_T-\gamma_0$, as $x\mapsto(2/x)^{1/3}$ is smooth near $f_0''(u_0)$. Finally $1+t_1u_+\ge c_0$ because $u_+\in J$.

(ii) Since $1+t_1u_0=1-t_1/(v(1+\sqrt r))\ge\eta$, the interval $J$ of (i) can be shrunk, uniformly in $r$, so that in addition $1+(1-\eta)v(1+\sqrt r)u\ge\eta/2$ on $J$; as $u<0$, then $1+t_1u\ge\eta/2$ on $J$, and $J\subset(-1/t_1,0)$ again. The term $m^{-1}t_1/(1+t_1u)$ of $f$, and the term $m^{-1}v/(1+vu)$ of $f_0$ it replaces, are bounded on $J$ with their first three derivatives by $C/m$. The argument of (i) applied to the indices $j\ge2$ then gives $\sup_J|f^{(l)}-f_0^{(l)}|\le C(|\mu|+s_2+m^{-1})$ for $l=0,1,2$, and the rest of the proof is unchanged, with $m\ge m_0$ ensuring the required smallness.
\end{proof}

\section{Proof of Theorem~\ref{thm:main}}\label{sec:proof}
\emph{A good event for the data.} Write $u_j=g_j/\|g_j\|$ with independent standard Gaussian $g_j\in\RR^{d}$; then $x_j=\sqrt d\,u_j$ has the required law. By Markov's inequality applied to $\exp(-\|g_j\|^2/2)$, $\PP(\|g_j\|^2<d/2)\le\exp(-d((\log2)/2-1/4))$, so $\min_j\|g_j\|^2\ge d/2$ with probability tending to one. For $G=(g_1,\dots,g_N)$, take $1/4$-nets of the unit spheres of $\RR^{d}$ and $\RR^{N}$, of sizes at most $9^{d}$ and $9^{N}$~\cite[Corollary~4.2.13]{Vershynin}. For fixed net points $x,y$, $x^{T}Gy$ is standard Gaussian, so the maximum of $|x^{T}Gy|$ over the nets is at most $4\sqrt{d+N}$ except with probability $2\exp(-(8-\log9)(d+N))$, and then $\|G\|_{\rm op}\le8\sqrt{d+N}$. Hence $\|U\|_{\rm op}\le\|G\|_{\rm op}/\min_j\|g_j\|\le C_U$ with probability tending to one, for a constant $C_U$ depending only on $\alpha$. Let
\[
\mathcal G_N=\Bigl\{\|U\|_{\rm op}\le C_U,\ \|B\|_{\rm op}\le N^{-1/8},\ \|R\|_F\le N^{-1/8},\ \Tr D^2\le N^{1/6}\Bigr\}.
\]
By Lemma~\ref{lem:cov} and Markov's inequality, the last three conditions fail with probabilities $O(N^{-1/2})$, $O(N^{-3/4})$ and $O(N^{-1/6})$, so $\PP(\mathcal G_N)\to1$.

\emph{The covariance on the good event.} On $\mathcal G_N$, \eqref{eq:decomp} and $N/d\to\alpha$ give $\|D\|_{\rm op}\le C$, and \eqref{eq:T} gives $\|T-vI_M\|_{\rm op}\le q/d+qN^{-1/8}+N^{-1/8}\to0$, uniformly. By~\eqref{eq:A0}, in the notation of Lemma~\ref{lem:edge}, $|\mu|\le\|D\|_{\rm op}/M\le C/M$ and $s_2\le\Tr D^2/M\le N^{1/6}/M$. Lemma~\ref{lem:edge}, with $r=M/m$, gives uniformly on $\mathcal G_N$
\begin{equation}\label{eq:shift}
m^{2/3}|E_T-E_0|\le CN^{-1/6}\to0,\qquad\gamma_T-\gamma_0\to0,\qquad1+t_1u_+\ge c_0.
\end{equation}
In particular $T$ is bounded above and below by positive constants uniformly on $\mathcal G_N$, and Proposition~\ref{prop:qf} gives, for each fixed $p$, the quadratic form bound with a constant independent of $N$ and of $U\in\mathcal G_N$.

\emph{Conditional edge universality.} Given $U$, the sample vectors $z_i=Q^{T}\sigma(U^{T}w_i)$, $i\le m$, are independent and centered with covariance $T$. Apply Corollary~1.4 of Chen and Lyu~\cite{CL}, with their sample size $N$ equal to our $m$, their dimension $M$ equal to our $M=N-1$, their vectors $q_i$ equal to our $z_i$ and their $\Sigma$ equal to $T$, to
\[
H_P=\frac1m\sum_{i=1}^{m}z_iz_i^{T}=\frac1m\,Q^{T}Y^{T}YQ.
\]
Let $U=U_N\in\mathcal G_N$ be a deterministic sequence. The vectors $z_i$ are independent and centered with common covariance $T$, which is Assumption~1.1(i) of~\cite{CL}; $M/m\to1/b\ne1$ is Assumption~1.1(ii); Proposition~\ref{prop:qf} with $C_T=\sup_N\|T\|_{\rm op}$ gives Assumption~1.1(iii) with $C_p=2\sqrt2\,C_U\kappa_4\,3^{k/2}\sqrt{C_T}\,(p-1)^{k}$; and the bounds on $T$ with~\eqref{eq:shift} give Assumption~1.2. Hence for every such sequence
\begin{equation}\label{eq:cond}
\gamma_Tm^{2/3}\bigl[\lambda_{\max}(H_P)-E_T\bigr]\Longrightarrow\TW.
\end{equation}
The convergence of the conditional distribution functions in~\eqref{eq:cond} is uniform over $U\in\mathcal G_N$: otherwise there would be $x\in\RR$, $\eta>0$ and $U_N\in\mathcal G_N$, along a subsequence, with conditional distribution function at $x$ at distance at least $\eta$ from that of $\TW$; completing the subsequence by arbitrary elements of $\mathcal G_N$ gives a deterministic sequence that contradicts~\eqref{eq:cond}. By~\eqref{eq:shift}, $E_T$ and $\gamma_T$ may be replaced by $E_0=v(1+\sqrt{M/m})^2$ and $\gamma_0$, uniformly on $\mathcal G_N$. Integrating over $U$ and using $\PP(\mathcal G_N)\to1$,
\begin{equation}\label{eq:uncond}
\gamma_0m^{2/3}\Bigl[\lambda_{\max}(H_P)-v\bigl(1+\sqrt{M/m}\bigr)^2\Bigr]\Longrightarrow\TW.
\end{equation}

\emph{Normalization.} The nonzero eigenvalues of $(YQ)^{T}(YQ)$ and $(YQ)(YQ)^{T}=YPY^{T}$ coincide, so $\lambda_{\max}(H_P)=(N/m)\lambda_{\max}(K_P)$. With $c_N=M/N=1-1/N$ and $M/m=c_N/b_N$, \eqref{eq:uncond} becomes
\[
\frac{b_N^{1/6}c_N^{1/6}}{v(\sqrt{b_N}+\sqrt{c_N})^{4/3}}\,N^{2/3}\Bigl[\lambda_{\max}(K_P)-v\bigl(\sqrt{b_N}+\sqrt{c_N}\bigr)^2\Bigr]\Longrightarrow\TW.
\]
Since $b_N$ stays in a compact subset of $(0,\infty)$, $v(\sqrt{b_N}+\sqrt{c_N})^2-v(1+\sqrt{b_N})^2=O(N^{-1})$, which is $o(N^{-2/3})$, and the ratio of the scaling parameter on the left to that in~\eqref{eq:main} tends to one. This proves~\eqref{eq:main}; tightness of the left side of~\eqref{eq:main} and $b_N\to b$ give the limit in probability. \qed

\begin{proof}[Proof of Corollary~\ref{cor:K}]
Repeat the argument with $Q=I_N$, $M=N$ and $z_i=\sigma(U^{T}w_i)$, whose covariance given $U$ is $T=\Sigma$; Proposition~\ref{prop:qf} applies with $Q=I_N$. Let $t_1\ge\dots\ge t_N$ be the eigenvalues of $\Sigma$ and $r=N/m$. By~\eqref{eq:decomp}, $\Sigma=\Sigma_0+(q/d)\mathbf 1\mathbf 1^{T}$ with $\Sigma_0=(v-q/d)I_N+qB+R$. On $\mathcal G_N$, $\|\Sigma_0-vI_N\|_{\rm op}\le q/d+qN^{-1/8}+N^{-1/8}$; by Weyl's inequality and the interlacing of eigenvalues under a positive semidefinite perturbation of rank one, $t_2,\dots,t_N$ lie within that distance of $v$, and $|t_1-v-qN/d|$ is at most that distance. Since $qN/d\to q\alpha<v\sqrt r$ in the limit, there is $\eta>0$ with $t_1\le(1-\eta)v(1+\sqrt r)$ on $\mathcal G_N$ for $N$ large, and $t_N\ge v/2$. Because $\Tr D=0$, $\sum_{j\ge2}(t_j-v)=-(t_1-v)$, so in the notation of Lemma~\ref{lem:edge}(ii) $|\mu|\le C/N$ and $s_2\le\Tr D^2/N\le N^{1/6}/N$. Lemma~\ref{lem:edge}(ii) then gives~\eqref{eq:shift} on $\mathcal G_N$, now with $E_0=v(1+\sqrt{N/m})^2$, together with $c_0\le\Sigma\le C$, which is Assumption~1.2 of~\cite{CL}. The rest of the argument is unchanged and gives~\eqref{eq:uncond} for $H=m^{-1}Y^{T}Y$ with $M=N$. The nonzero eigenvalues of $H$ and $(N/m)K$ coincide, so $\lambda_{\max}(H)=(N/m)\lambda_{\max}(K)$, and the normalization step with $c_N=1$ gives~\eqref{eq:main} for $K$.
\end{proof}

\begin{remark}[Finite $N$]\label{rem:finite}
Given $W$, the columns of $Y$ are independent and identically distributed in $\RR^{m}$; for $\sigma=h_2$ their covariance is $\Sigma_W=\frac{d}{d+2}\,G_W\circ G_W-\frac1{d+2}\,s_Ws_W^{T}$, where $G_W=WW^{T}/d$, $s_W$ is the diagonal of $G_W$ and $\circ$ is the entrywise product. Let $E_W$ be $(N-1)/N$ times the right edge $E_T$ of Section~\ref{sec:edge} for the covariance $T=\Sigma_W$, with the roles of the dimension and the sample size exchanged: the dimension is $m$ and there are $N-1$ samples. Heuristically $\lambda_{\max}(K_P)$ lies near $E_W$; the $N-1$ columns of $YQ$ are uncorrelated with covariance $\Sigma_W$ but not independent, and the location of $\lambda_{\max}(K_P)$ relative to $E_W$ is open. The diagonal of $\Sigma_W$ carries the spread of $|w_i|^2/d$, of variance $2/d$, which displaces $E_W$ from $E_*=(\sqrt{b_N}+\sqrt{c_N})^2$, the centering of the proof above for $v=1$, with $c_N=1-1/N$. For $m=N/2$ and $d=3N/4$, the mean of $(c_Nb_N)^{1/6}(\sqrt{b_N}+\sqrt{c_N})^{-4/3}N^{2/3}(E_W-E_*)$ over six independent $W$ is $2.4$, $2.1$, $1.5$, $1.4$ and $1.0$ at $N=500$, $1000$, $2000$, $4000$ and $8000$, which agrees with a displacement of order $N^{-1}$, that is, the rate $N^{-1/3}$ in these units; for $m=2N$ and $d=3N/2$ it is $0.6$ and $0.7$ at $N=1000$ and $2000$. A displacement of order $N^{-1}$ is negligible on the scale $N^{-2/3}$ but not for $N$ in the thousands.
\end{remark}

\section*{Acknowledgements and funding}
This work was done by multiple AI agents, Claude (Anthropic) and Codex (OpenAI), connected to Math Brain~\cite{MathBrainRepo} and working together inside the Brain framework~\cite{BrainFramework}, a shared, persistent memory in which the agents record claims, proofs, failed attempts and methods, retrieve one another's results, coordinate who works on what, and audit each other's work. For more on the details behind the Brain and Math Brain, see the Math Brain paper~\cite{MathBrainPaper}, the Math Brain repository~\cite{MathBrainRepo} and the Brain framework repository~\cite{BrainFramework}.

The research presented here was supported by the European Research Council (ERC) under the European Union's Horizon Europe research and innovation program (grant agreement No.~101041711) and by the Simons Foundation as part of the Collaboration on the Mathematical and Scientific Foundations of Deep Learning. It was also supported by Heights Labs, by Convex Nexus Capital and by the Israel Science Foundation (grants 2258/19 and 4101/25).

\section*{Declaration of competing interest}
The author holds equity in Heights Labs and in Convex Nexus Capital and has a management role at Convex Nexus Capital. These entities are included in the funding declaration above.

\section*{Declaration of generative AI and AI-assisted technologies}
Codex (OpenAI, GPT-6.1) developed the proofs and wrote the first verification script; Claude Opus~5.5 (Anthropic) checked the proofs, wrote the second and third scripts and wrote the paper. The two systems worked through the shared memory of Math Brain~\cite{MathBrainPaper,MathBrainRepo} and the Brain framework~\cite{BrainFramework}. The results rest on the edge universality theorem of Chen and Lyu~\cite{CL}, the quadratic form bound of Fan, Ma, Paquette and Wang~\cite{FMPW} and the edge function of Silverstein and Choi~\cite{SC}; the covariance after the row means are removed and the displacement of the right edge (Lemmas~\ref{lem:cov}, \ref{lem:T} and~\ref{lem:edge}) are new. The author directed the work and is responsible for the content of the paper and for decisions concerning its dissemination.

\section*{Data availability}
The \LaTeX{} source of the paper and three Python scripts are available at \url{https://github.com/yspennstate/row-centered-conjugate-kernel-tracy-widom/tree/v1.6}. The first evaluates, in exact rational arithmetic, both sides of each inequality in the proof of Proposition~\ref{prop:qf} for $\sigma$ of degree two to five and several matrices $U$. The second evaluates~\eqref{eq:sphmom}, the closed form and the bound for $\mu_4$ in Lemma~\ref{lem:cov}, the constants of Section~\ref{sec:edge} and the change of normalization in Section~\ref{sec:proof}. The third computes the edges $E_W$ quoted in Remark~\ref{rem:finite}.

\appendix
\section{Notation}\label{app:notation}
Symbols are listed in the order in which they first appear. Standard symbols ($\RR$, $\EE$, $\PP$, $\Tr$, $\Var$) are omitted, and so are generic constants $C$, $c$ and letters defined and used inside a single proof.

\begin{center}\small\begin{tabular}{@{}lp{0.6\textwidth}l@{}}\toprule
Symbol & Meaning & First appears \\ \midrule
$m,d,N$ & number of hidden units, input dimension, number of data points & Section~\ref{sec:intro} \\
$W$ & $m\times d$ weight matrix with independent standard Gaussian entries & Section~\ref{sec:intro} \\
$X$, $x_j$ & $d\times N$ data matrix and its columns, uniform on the sphere of radius $\sqrt d$ & Section~\ref{sec:intro} \\
$\sigma$, $k$ & activation, a polynomial of degree $k\ge2$ & Section~\ref{sec:intro} \\
$Y$ & the feature matrix $\sigma(WX/\sqrt d)$ & Section~\ref{sec:intro} \\
$K$ & conjugate kernel $YY^{T}/N$ & Section~\ref{sec:intro} \\
$Z$ & a standard Gaussian variable & Section~\ref{sec:intro} \\
$\alpha$ & limit of $N/d$ & Section~\ref{sec:intro} \\
$b_N$, $b$ & the ratio $m/N$ and its limit & Section~\ref{sec:intro} \\
$\mathbf 1$, $e$ & vector of ones in $\RR^N$, and $e=N^{-1/2}\mathbf 1$ & Section~\ref{sec:intro} \\
$P$ & projection $I_N-\mathbf 1\mathbf 1^{T}/N$ removing the row means & Section~\ref{sec:intro} \\
$K_P$ & row-centered conjugate kernel $YPY^{T}/N$ & Section~\ref{sec:intro} \\
$v$ & $\EE\sigma(Z)^2$ & Section~\ref{sec:intro} \\
$\TW$ & Tracy--Widom law of the Gaussian orthogonal ensemble & Section~\ref{sec:intro} \\
$\lambda_{\max}$, $\lambda_2$ & largest and second largest eigenvalue & Section~\ref{sec:intro} \\
$h_j$, $a_j$ & normalized Hermite polynomials $\mathrm{He}_j/\sqrt{j!}$ and the coefficients of $\sigma$ on them & Section~\ref{sec:intro} \\
$q$ & $a_2^2$ & Section~\ref{sec:intro} \\
$g$ & standard Gaussian vector in $\RR^d$ & Section~\ref{sec:intro} \\
$U$, $u_j$ & the matrix with unit columns $u_j=x_j/\sqrt d$ & Section~\ref{sec:intro} \\
$z_i$, $\xi_i$ & sample vectors of a sample covariance matrix, and $T^{-1/2}z_i$ & Section~\ref{sec:intro} \\
$M$ & dimension of the sample vectors, equal to $N-1$ in the proofs & Section~\ref{sec:intro} \\
$T$ & covariance of the sample vectors & Section~\ref{sec:intro} \\
$E_T$, $\gamma_T$ & right edge determined by $T$ and $M/m$, and its scaling parameter & Section~\ref{sec:intro} \\
$t_1\ge\dots\ge t_M$ & eigenvalues of $T$; $t_{\max}=t_1$ & Section~\ref{sec:intro} \\
$u_+$ & the zero of $f'$ defining $E_T$ & Section~\ref{sec:intro} \\
$\tau$ & constant in the regularity condition of the right edge & Section~\ref{sec:intro} \\
$w$ & weight vector of the random features model $\sigma(Xw)$ of Fan, Ma, Paquette and Wang & Section~\ref{sec:intro} \\
$\|\cdot\|_p$ & norm of $L^p$, $(\EE|V|^p)^{1/p}$ & Section~\ref{sec:intro} \\
$\|\cdot\|_F$, $\|\cdot\|_{\rm op}$ & Hilbert--Schmidt norm and operator norm & Section~\ref{sec:intro} \\
$\Sigma$ & covariance of a row of $Y$ given the data & Section~\ref{sec:intro} \\
$B$, $R$ & zero-diagonal matrices with $B_{jl}=\rho_{jl}^2-1/d$ and $R_{jl}$ the terms of $\Sigma_{jl}$ of degree at least three in $\rho_{jl}$ & Section~\ref{sec:intro} \\
$Q$ & $N\times M$ matrix with orthonormal columns and $QQ^{T}=P$ & Section~\ref{sec:intro} \\
$A_0$ & $T-vI_M$ & Section~\ref{sec:intro} \\
$O_P$, $o_P$ & bounded in probability, and tending to zero in probability & Section~\ref{sec:intro} \\
$\bar y$ & vector of row means $Y\mathbf 1/N$ & Remark~\ref{rem:scope} \\
$\phi$ & $\sigma(U^{T}g)$, a row of $Y$ given the data & Section~\ref{sec:qf} \\
$C_U$, $C_T$ & bounds on $\|U\|_{\rm op}$ and $\|T\|_{\rm op}$ & Section~\ref{sec:qf} \\
$D_4$ & $(\EE|\sigma'(Z)|^4)^{1/4}$ & Section~\ref{sec:qf} \\
$\rho_{jl}$ & $u_j^{T}u_l$ & Section~\ref{sec:cov} \\
$D$ & $\Sigma-vI_N$ & Section~\ref{sec:cov} \\
$f$ & edge function of Silverstein and Choi & Section~\ref{sec:edge} \\
$r$ & $M/m$ & Section~\ref{sec:edge} \\
$u_0$, $E_0$, $\gamma_0$ & $u_+$, $E_T$, $\gamma_T$ for $T=vI_M$ & Section~\ref{sec:edge} \\
$\mu$, $s_2$ & mean and mean square of $t_j-v$ & Section~\ref{sec:edge} \\
$\mathcal G_N$ & event of probability tending to one on which the data are good & Section~\ref{sec:proof} \\
$H_P$ & $m^{-1}\sum_iz_iz_i^{T}$ for $z_i=Q^{T}\phi_i$ & Section~\ref{sec:proof} \\ \bottomrule
\end{tabular}\end{center}


\begin{thebibliography}{99}
\bibitem{BBP} J.~Baik, G.~Ben~Arous and S.~P\'ech\'e, \emph{Phase transition of the largest eigenvalue for nonnull complex sample covariance matrices}, Ann. Probab. 33 (2005), 1643--1697, arXiv:math/0403022.
\bibitem{BPZ} Z.~Bao, G.~Pan and W.~Zhou, \emph{Universality for the largest eigenvalue of sample covariance matrices with general population}, Ann. Statist. 43 (2015), 382--421, arXiv:1304.5690.
\bibitem{BP21} L.~Benigni and S.~P\'eché, \emph{Eigenvalue distribution of some nonlinear models of random matrices}, Electron. J. Probab. 26 (2021), Paper No.~150, arXiv:1904.03090.
\bibitem{BP22} L.~Benigni and S.~P\'eché, \emph{Largest eigenvalues of the conjugate kernel of single-layered neural networks}, preprint (2022), arXiv:2201.04753.
\bibitem{CL} G.~Chen and X.~Lyu, \emph{Tracy--Widom law for sample covariance matrices under quadratic form concentration}, preprint (2026), arXiv:2610.41416v2.
\bibitem{Chernoff} H.~Chernoff, \emph{A note on an inequality involving the normal distribution}, Ann. Probab. 9 (1981), 533--535.
\bibitem{EK10} N.~El~Karoui, \emph{The spectrum of kernel random matrices}, Ann. Statist. 38 (2010), 1--50, doi:10.1214/08-AOS648.
\bibitem{ElKaroui} N.~El~Karoui, \emph{Tracy--Widom limit for the largest eigenvalue of a large class of complex sample covariance matrices}, Ann. Probab. 35 (2007), 663--714, arXiv:math/0503109.
\bibitem{FM} Z.~Fan and A.~Montanari, \emph{The spectral norm of random inner-product kernel matrices}, Probab. Theory Related Fields 173 (2019), 27--85, doi:10.1007/s00440-018-0830-4.
\bibitem{FP} D.~F\'eral and S.~P\'ech\'e, \emph{The largest eigenvalues of sample covariance matrices for a spiked population: diagonal case}, J. Math. Phys. 50 (2009), 073302, arXiv:0812.2320.
\bibitem{FMPW} Z.~Fan, R.~Ma, E.~Paquette and Z.~Wang, \emph{Anisotropic local law for non-separable sample covariance matrices}, preprint (2026), arXiv:2602.17960v2.
\bibitem{FW} Z.~Fan and Z.~Wang, \emph{Spectra of the conjugate kernel and neural tangent kernel for linear-width neural networks}, Advances in Neural Information Processing Systems 33 (2020), arXiv:2005.11879.
\bibitem{Hariya} Y.~Hariya, \emph{A unification of hypercontractivities of the Ornstein--Uhlenbeck semigroup and its connection with $\Phi$-entropy inequalities}, J. Funct. Anal. 275 (2018), 2647--2683, doi:10.1016/j.jfa.2018.08.009, arXiv:1707.03163.
\bibitem{Johnstone} I.~M.~Johnstone, \emph{On the distribution of the largest eigenvalue in principal components analysis}, Ann. Statist. 29 (2001), 295--327.
\bibitem{LS} J.~O.~Lee and K.~Schnelli, \emph{Tracy--Widom distribution for the largest eigenvalue of real sample covariance matrices with general population}, Ann. Appl. Probab. 26 (2016), 3786--3839, arXiv:1409.4979.
\bibitem{LLC} C.~Louart, Z.~Liao and R.~Couillet, \emph{A random matrix approach to neural networks}, Ann. Appl. Probab. 28 (2018), 1190--1248, arXiv:1702.05419.
\bibitem{MM} R.~Ma and T.~Misiakiewicz, \emph{The anisotropic local law for sample covariance matrices under quadratic-form concentration}, preprint (2026), arXiv:2609.09440v2.
\bibitem{MP} V.~A.~Marchenko and L.~A.~Pastur, \emph{Distribution of eigenvalues for some sets of random matrices}, Mat. Sb. 72(114) (1967), 507--536.
\bibitem{Nelson} E.~Nelson, \emph{The free Markoff field}, J. Funct. Anal. 12 (1973), 211--227.
\bibitem{ODonnell} R.~O'Donnell, \emph{Analysis of Boolean functions}, Cambridge University Press, Cambridge, 2014.
\bibitem{PW} J.~Pennington and P.~Worah, \emph{Nonlinear random matrix theory for deep learning}, Advances in Neural Information Processing Systems 30 (2017); J. Stat. Mech. Theory Exp. (2019), 124005, doi:10.1088/1742-5468/ab3bc3.
\bibitem{SC} J.~W.~Silverstein and S.-I.~Choi, \emph{Analysis of the limiting spectral distribution of large dimensional random matrices}, J. Multivariate Anal. 54 (1995), 295--309.
\bibitem{TW} C.~A.~Tracy and H.~Widom, \emph{On orthogonal and symplectic matrix ensembles}, Comm. Math. Phys. 177 (1996), 727--754.

\bibitem{TW94} C.~A.~Tracy and H.~Widom, \emph{Level-spacing distributions and the Airy kernel}, Comm. Math. Phys. 159 (1994), 151--174.
\bibitem{Vershynin} R.~Vershynin, \emph{High-dimensional probability}, Cambridge University Press, Cambridge, 2018.
\bibitem{WWF} Z.~Wang, D.~Wu and Z.~Fan, \emph{Nonlinear spiked covariance matrices and signal propagation in deep neural networks}, Proceedings of Machine Learning Research 247 (2024), arXiv:2402.10127.
\bibitem{MathBrainRepo} Y.~Shmalo, \emph{Math Brain: public snapshot, code and tests}, 2026, \url{https://github.com/yspennstate/math-brain-public/tree/8c9a6bd478c4ec4082347151b7ae16db50c917dd}.
\bibitem{MathBrainPaper} Y.~Shmalo, \emph{Math Brain: shared, auditable memory for multi-agent mathematical research}, October 2026, doi:10.5281/zenodo.24150837, \url{https://github.com/yspennstate/math-brain-public/tree/8c9a6bd478c4ec4082347151b7ae16db50c917dd/papers/math-brain-2026-10-24}.
\bibitem{BrainFramework} Y.~Shmalo, \emph{Brain Framework}, public software repository, 2026, \url{https://github.com/yspennstate/brain-framework-public/tree/5629e644cdd510173ae97e9aa92642dabd22dc57}.
\end{thebibliography}
\end{document}
